OpenFisca-France 176.1 (législation des revenus 2024) ; calculs de l'auteur.\par Note : salariés, en \% du salaire brut ; IR, CSG et CRDS ; le couple partage ses revenus 60/40. Foyers sans aucune niche ni déduction : un foyer réel de même revenu paie souvent moins aujourd'hui. Barème proposé, sur le revenu après cotisations, par part : taux marginal de 7 \% au premier euro, 14 \% à 20 000 \euro{}, 21,3 \% à 30 000 \euro{}, 45 \% à partir de 120 000 \euro{} ; deux parts pour un couple ; crédit de 750 \euro{} par enfant (enfants de plus de trois ans ici).
